\documentclass[british]{article}

\usepackage[margin=1.5in]{geometry}

\title{Serre's theorem}
\author{P. Gabriel}
\date{1958}

% Standard
\usepackage{amssymb,amsmath,amscd}
\usepackage{enumerate}
\usepackage{tikz-cd}
\usepackage{mathtools}
\usepackage{booktabs}


%Bibliography
\usepackage[backend=biber,maxnames=99,maxalphanames=4]{biblatex}
\addbibresource{SCC-4-2.bib}


% Colours
\usepackage{xcolor}


% Fonts
\usepackage{mathrsfs}
\usepackage{fouriernc}


% Things for Quarto/pandoc
\definecolor{quarto-callout-color-frame}{HTML}{000000}
\usepackage{calc,array}
\usepackage{longtable}
\providecommand{\tightlist}{%
\setlength{\itemsep}{0pt}\setlength{\parskip}{0pt}}
\newcounter{none}
\usepackage{graphicx}
\makeatletter
\newsavebox\pandoc@box
\newcommand*\pandocbounded[1]{% scales image to fit in text height/width
  \sbox\pandoc@box{#1}%
  \Gscale@div\@tempa{\textheight}{\dimexpr\ht\pandoc@box+\dp\pandoc@box\relax}%
  \Gscale@div\@tempb{\linewidth}{\wd\pandoc@box}%
  \ifdim\@tempb\p@<\@tempa\p@\let\@tempa\@tempb\fi% select the smaller of both
  \ifdim\@tempa\p@<\p@\scalebox{\@tempa}{\usebox\pandoc@box}%
  \else\usebox{\pandoc@box}%
  \fi%
}
% Set default figure placement to htbp
\def\fps@figure{htbp}
\makeatother
\usepackage{kvsetkeys}


% Header and footer
\usepackage{fancyhdr}
\usepackage{lastpage}
\usepackage{datetime2}
\pagestyle{fancy}
\fancypagestyle{plain}{}
\fancyhf{}
\renewcommand{\headrulewidth}{0pt}
\cfoot{\small\thepage\ of \pageref*{LastPage}}
\lfoot{\footnotesize Build date: \today}


% Theorem environments
\usepackage{amsthm}
\newenvironment{itenv}[1]
  {\phantomsection\par\smallskip\noindent\textbf{#1.}\itshape}
  {\par\smallskip}
\newenvironment{rmenv}[1]
  {\phantomsection\par\smallskip\noindent\textbf{#1.}\rmfamily}
  {\par\smallskip}


% Shortcuts
\newcommand{\oldpage}[1]{\marginpar{\footnotesize$\Big\vert$ \textit{p.~#1}}}


%% Document %%

% hyperref last, as is "good practice"

\usepackage{hyperref}
\hypersetup{colorlinks,linkcolor={purple!50!black},citecolor={purple!50!black},urlcolor={purple!80!black}}

\usepackage{embedall}
\begin{document}

\maketitle
\thispagestyle{fancy}

\setcounter{footnote}{0}


%% Content %%

\emph{This document is a translation into English of the following:}

\begin{quote}
Gabriel, P. ``Le théorème de Serre''. \emph{Séminaire Claude Chevalley}
\textbf{4} (1958--59), Talk no. 2.
\href{http://www.numdam.org/item/SCC_1958-1959__4__A2_0}{numdam.org/item/SCC\_1958-1959\_\_4\_\_A2\_0}
\end{quote}

\emph{The translator (\href{https://thosgood.net}{Tim Hosgood}) takes
full responsibility for any errors introduced in this document, and
claims no rights to any of the mathematical content.}

\tableofcontents

\providecommand{\scr}[1]{{\mathscr{#1}}}
\renewcommand{\cal}[1]{{\mathcal{#1}}}
\renewcommand{\frak}[1]{{\mathfrak{#1}}}
\renewcommand{\geq}{\geqslant}
\renewcommand{\leq}{\leqslant}

\section*{Affine algebraic sets and classical
functors}\label{affine-algebraic-sets-and-classical-functors}
\addcontentsline{toc}{section}{Affine algebraic sets and classical
functors}

\section{Ringed topological spaces}\label{section-1}

\oldpage{2-01}

From now on, and unless otherwise mentioned, the rings we consider will
be assumed to be commutative, with a unit element, and
\emph{Noetherian}. We define a \emph{ringed topological space}
\((V,{\mathscr{A}})\) to be a topological space \(V\) endowed with the
structure defined by the data of a sheaf of rings \({\mathscr{A}}\). If
\((V,{\mathscr{A}})\) and \((W,{\mathscr{B}})\) are two ringed
topological spaces, then we define a morphism from the first to the
second by the data of:

\begin{enumerate}
\def\labelenumi{\alph{enumi}.}
\tightlist
\item
  a continuous map \(\psi\colon V\to W\);
\item
  for every open subset \(U\) of \(W\), a ring homomorphism \[
     \varphi_U\colon{\mathscr{B}}(U)\to{\mathscr{A}}(\varphi^{-1}(U))
   \] that is compatible with the restriction maps.
\end{enumerate}

The composition of two morphisms is defined in the evident way, and we
speak of the category of ringed topological spaces. In what follows,
\(V\) will almost always be a Zariski space. To such a \(V\), we often
associate a topological space \(S(V)\), which is \emph{the scheme of
\(V\)}, and which is defined in the following way:

\begin{itemize}
\tightlist
\item
  the points of \(S(V)\) are the closed irreducible subsets of \(V\);
\item
  the closed subsets of \(S(V)\) are the sets \({\mathscr{F}}\), where
  \(F\) is a closed subset of \(V\), and \({\mathscr{F}}\) denotes the
  set of closed irreducible subsets of \(V\) that are contained in
  \(F\).
\end{itemize}

It is then clear that the correspondence \(F\mapsto{\mathscr{F}}\)
between closed subsets of \(V\) and closed subsets of \(S(V)\) is
bijective, that \(S(V)\) is a Zariski space, and that every closed
irreducible subset of \(S(V)\) is the closure of a unique point. To
every sheaf on \(V\), we canonically associate a sheaf on \(S(V)\). In
particular, if \((V,{\mathscr{A}})\) is a ringed topological space, and
if \(V\) is a Zariski space, then we denote by
\((S(V),S({\mathscr{A}}))\) the \emph{scheme of \((V,{\mathscr{A}})\)},
which is defined to be the ringed topological space given by \(S(V)\)
and the sheaf associated to \({\mathscr{A}}\). \oldpage{2-02} Of course,
the scheme of \((S(V),S({\mathscr{A}}))\) is isomorphic to
\((S(V),S({\mathscr{A}}))\).

If \(A\) is a Noetherian Jacobson ring, and
\((\Omega(A),{\mathscr{U}})\) is the ringed topological space defined by
its maximal spectrum, then the associated scheme is exactly
\((V(A),{\mathscr{U}})\) (see Talk no. 1). If \((V,{\mathscr{A}})\) is
an algebraic set endowed with the sheaf of germs of regular functions,
then the associated scheme has been defined by Chevalley in \autocite{3}
and \autocite{4}. In this case, \(V\) is the subspace of \(S(V)\) given
by the closed points, and \({\mathscr{A}}\) is the restriction of
\(S({\mathscr{A}})\) to \(V\). From now on, we will almost always
restrict to the case of algebraic sets over an algebraically closed
field.

\section{Sheaves of ideals of an algebraic set}\label{section-2}

If \(V\) is an algebraic set, and \({\mathscr{A}}\) its sheaf of germs
of regular functions, then the notion of quasi-coherent (resp. coherent)
sheaves generalises the notion of a module (resp. module of finite type)
over the ring of coordinates of an affine algebraic set. We similarly
generalise the notions of support and dimension: if \({\mathscr{F}}\) is
a coherent algebraic sheaf on \(V\), then its \emph{support} is the set
of points \(x\) of \(V\) where the fibre \({\mathscr{F}}_x\) of
\({\mathscr{F}}\) is not zero (if the fibre of \({\mathscr{F}}\) is zero
at \(x\), then it is zero on a neighbourhood of \(x\), since
\({\mathscr{F}}\) is coherent, and thus the support is closed). We then
define the \emph{dimension of \({\mathscr{F}}\)} to be the dimension of
its support. When \(V\) is affine and the coherent sheaf
\({\mathscr{F}}\) is associated to a module \(M\) of finite type, then
the support of \({\mathscr{F}}\) is given by the set of maximal ideals
of the coordinate ring that contain the annihilator of \(M\).

The notion of an ideal generalises in the following way: we define a
\emph{sheaf of ideals} of \((V,{\mathscr{A}})\) to be any quasi-coherent
(and thus coherent) subsheaf of \({\mathscr{A}}\). We then find the
inevitable correspondence between ideals and closed subsets:

If \({\mathscr{I}}\) is a sheaf of ideals, then we associate to it the
following closed subset \(W({\mathscr{I}})\) of \(V\):

\begin{quote}
\(x\in W({\mathscr{I}})\) if and only if \({\mathscr{I}}_x\), the fibre
of \({\mathscr{I}}\) at \(x\), is a proper ideal of \({\mathscr{O}}_x\)
(i.e.~\({\mathscr{I}}_x\neq{\mathscr{O}}_x\)), or if and only if the
germs of the functions defined by \({\mathscr{I}}\) at \(x\) are zero.
\(W({\mathscr{I}})\) is thus the support of
\({\mathscr{A}}|{\mathscr{I}}\).
\end{quote}

Conversely, if \(W\) is a closed subset of \(V\), then we associate to
it the following sheaf of ideals \({\mathscr{I}}(W)\): \oldpage{2-03}

\begin{quote}
The sections of \({\mathscr{I}}(W)\) on an open subset \(U\) are the
regular functions defined on \(U\) that are zero on \(U\cap W\). If
\({\mathscr{I}}(U,x)\) denotes the inverse image in \({\mathscr{A}}(U)\)
of the maximal ideal of \({\mathscr{O}}_x\), then
\(\Gamma(U,{\mathscr{I}}(W))\) is the intersection of the
\({\mathscr{I}}(U,x)\) where \(x\) runs over \(U\cap W\).
\end{quote}

If \(U\) is an affine open subset, then the sheaf of ideals
\({\mathscr{I}}(W)|U\) is exactly the sheaf associated to the ideal
\(\Gamma(U,{\mathscr{I}}(W))\) of \({\mathscr{A}}(U)\), which shows that
\({\mathscr{I}}(W)\) is coherent.

\protect\phantomsection\label{proposition-1}
\begin{itenv}{Proposition 1}

---

\begin{enumerate}
\def\labelenumi{\alph{enumi}.}
\tightlist
\item
  The map \(W\mapsto{\mathscr{I}}(W)\) gives a bijective correspondence
  between the closed subsets of \(V\) and the sheaves of ideals
  \({\mathscr{I}}\) that satisfy the following condition: for every
  \(x\in V\), \({\mathscr{I}}_x\) is either equal to \({\mathscr{O}}_x\)
  or equal to an intersection of prime ideals of \({\mathscr{O}}_x\).
\item
  The map \(W\mapsto{\mathscr{I}}(W)\) associates the closed subsets
  whose connected components are all irreducible to the sheaves of prime
  ideals (for every \(x\), \({\mathscr{I}}_x\) is either equal to
  \({\mathscr{O}}_x\) or equal to a prime ideal).
\end{enumerate}

\end{itenv}

\begin{proof}
---

\begin{enumerate}
\def\labelenumi{\alph{enumi}.}
\tightlist
\item
  It suffices to give a proof in the case where \((V,{\mathscr{A}})\) is
  an affine algebraic set. So let \(A={\mathscr{A}}(V)\), and
  \(\mathfrak{a}=\Gamma(V,{\mathscr{I}}(W))\). We have already seen that
  \(\mathfrak{a}\) is then the intersection of the prime ideals
  \({\mathscr{I}}(V,x)\), where \(x\) runs over \(W\). So
  \(\mathfrak{a}\) is an intersection of prime ideals. Since the
  correspondence between intersections of prime ideals of \(A\) and
  closed subsets of \(V\) is bijective, so too is the correspondence
  between ideals of \(A\) and sheaves of ideals of
  \((V,{\mathscr{A}})\); it remains only to show, conversely, that the
  sheaf associated to an intersection \(\mathfrak{a}\) of prime ideals
  satisfies the condition of the proposition: this follows from the
  conservation properties of the prime decomposition under localisation.
\item
  The proof is analogous.
\end{enumerate}

\end{proof}

\section{The ring of rational functions of an algebraic
set}\label{section-3}

We recall that, if \(V\) is an algebraic set, then we define a
\emph{rational function on \(V\)} to be a regular map \(f\) from an
everywhere-dense open subset of \(V\) to the field of constants \(K\).
We further suppose that the domain of definition of \(f\) cannot be
extended. The rational functions on \(V\) form an algebra \(K(V)\) over
\(K\). If we denote by \(V_i\) the irreducible components of \(V\), then
\(K(V)\) is isomorphic to the product \(\prod_i K(V_i)\), with the
isomorphism being the obvious one. Finally, if \(V\) is irreducible, and
if \(U\) is an affine open subset of \(V\), then \(K(V)\) is exactly the
field of fractions of the coordinate ring of \(U\).

The sheaf \({\mathscr{K}}(V)\) of rational functions on \(V\) is then
defined in the following way:

\begin{quote}
to every open subset \(U\) of \(V\), we associate the ring \(K(U)\) of
rational functions on \(U\), with the restrictions being obvious.
\end{quote}

\oldpage{2-04}

It is clear that this defines a quasi-coherent sheaf on \(V\).
Furthermore, if we denote by \(V_i\) the irreducible components of
\(V\), then \(K(U)\) is exactly the product
\(\prod_{V_i\cap U\neq\varnothing}K(V_i)\). It thus follows that the
sheaf \({\mathscr{K}}(V)\) is obtained in the following way: let
\({\mathscr{K}}_i\) be the sheaf on \(V\) that is zero outside of
\(V_i\), and that is constant with fibre \(K(V_i)\) on \(V_i\). Then
\({\mathscr{K}}(V)\) is the product of the sheaves \({\mathscr{K}}_i\).

\section{Characterisation of affine algebraic sets}\label{section-4}

Let \((V,{\mathscr{A}})\) be an algebraic set endowed with its sheaf of
rings, and let \(A=\Gamma(V,{\mathscr{A}})\) be the coordinate ring of
\(V\). We know that we then have a canonical morphism
\((V,{\mathscr{A}})\to(S(V),S({\mathscr{A}}))\), and that
\((V,{\mathscr{A}})\) is affine if and only if
\((S(V),S({\mathscr{A}}))\) is the prime spectrum of an algebra over
\(K\), of finite type and with no nilpotent elements.

We will now define a morphism
\((S(V),S({\mathscr{A}}))\to(V(A),{\mathscr{U}})\). For this, let \(x\)
be an arbitrary element of \(S(V)\), and let \(S({\mathscr{A}})_x\) be
the fibre of \(S({\mathscr{A}})\) at \(x\): this an a local ring, and we
have a restriction homomorphism \(A\to S({\mathscr{A}})_x\). We denote
by \(\mathfrak{p}_x\) the prime ideal of \(A\) given by the inverse
image under this homomorphism of the maximal ideal of
\(S({\mathscr{A}})_x\). It is an ideal of functions on \(V\) that are
zero on the closed point \(x\).

The map \(\varphi\colon x\mapsto\mathfrak{p}_x\) is a continuous map
from \(S(V)\) to \(V(A)\). Indeed, it suffices to show that the inverse
image of a special open subset \(U_f\) of \(V(A)\) is an open subset of
\(S(V)\). But \(\varphi^{-1}(U_f)\) consists of points \(x\) of \(S(V)\)
such that the image \(f_x\) of \(f\) in \({\mathscr{A}}_x\) is
invertible, and if \(f_x\) is invertible at \(x\), then \(f_y\) is
invertible for \(y\) in a neighbourhood of \(x\). QED.

It follows from the above that \(1/f\) is a section of
\(S({\mathscr{A}})\) over \(\varphi^{-1}(U_f)\), and so the canonical
map from \(A\) to \(\Gamma(\varphi^{-1}(U_f),S({\mathscr{A}}))\) extends
to a homomorphism: \[
  \varphi_{U_f}\colon A_f \to \Gamma(\varphi^{-1}(U_f), S({\mathscr{A}})).
\]

The \(\varphi_{U_f}\) are compatible with the restriction maps, and, by
``passing to the inductive limit'', they define a morphism of ringed
topological spaces: \oldpage{2-05} \[
  \varphi\colon (S(V),S({\mathscr{A}})) \to (V(A),{\mathscr{U}}).
\]

It is clear that \(\varphi\) is an isomorphism if and only if
\((V,{\mathscr{A}})\) is an affine algebraic set. The following theorem
further examines this case:

\begin{itenv}{Serre's Theorem}

The following are equivalent:

\begin{enumerate}
\def\labelenumi{\alph{enumi}.}
\tightlist
\item
  \((V,{\mathscr{A}})\) is an affine algebraic set.
\item
  If \(0\to{\mathscr{F}}\to{\mathscr{G}}\to{\mathscr{H}}\to0\) is an
  exact sequence of quasi-coherent sheaves, then the sections over \(V\)
  form an exact sequence.
\item
  If \(0\to{\mathscr{F}}\to{\mathscr{G}}\to{\mathscr{H}}\to0\) is an
  exact sequence of coherent sheaves, then the sections over \(V\) form
  an exact sequence.
\item
  There exist sections \(f_i\) of \({\mathscr{A}}\) over \(V\) such
  that:

  \begin{itemize}
  \tightlist
  \item
    the ideal generated by the \(f_i\) is equal to \(A\); and
  \item
    the open subsets \(V_{f_i}\) of \(V\), where the \(f_i\) are
    non-zero, are affine open subsets, and they cover \(V\).
  \end{itemize}
\end{enumerate}

\end{itenv}

\begin{proof}
It remains only to show that (c) implies (d), and that (d) implies (a).

\bigskip

For (c)\(\implies\)(d), consider:

\protect\phantomsection\label{lemma-1}
\begin{itenv}{Lemma 1}
If \({\mathscr{F}}\) is a non-zero coherent sheaf, then
\(\Gamma(V,{\mathscr{F}})\) is non-zero.

\end{itenv}

\begin{proof}
Indeed, the support of \({\mathscr{F}}\) contains a closed point, say
\(x\).

The fibre \({\mathscr{F}}_x\) of \({\mathscr{F}}\) is non-zero at \(x\),
and, if \(\mathfrak{m}_x\) denotes the maximal ideal of
\({\mathscr{A}}_x\), then
\({\mathscr{F}}_x|\mathfrak{m}_x{\mathscr{F}}_x\) is non-zero (by the
Nakayama lemma). The sheaf \({\mathscr{G}}\) that is zero away from
\(x\) and ``has the value''
\({\mathscr{F}}_x|\mathfrak{m}_x{\mathscr{F}}_x\) at \(x\) is then
coherent, and we clearly have a surjection: \[
  {\mathscr{F}}\to{\mathscr{G}}\to0.
\]

It thus follows that we have an epimorphism from
\(\Gamma(V,{\mathscr{F}})\) to
\(\Gamma(V,{\mathscr{G}}) = {\mathscr{F}}_x|\mathfrak{m}_x{\mathscr{F}}_x\),
and \(\Gamma(V,{\mathscr{F}})\) is non-zero.
\end{proof}

With this proven, let \(f\) be a section of \({\mathscr{A}}\) over \(V\)
(a regular function on \(V\)). If \(U\) is an arbitrary open subset of
\(V\), then we denote by \(U_f\) the set of points of \(U\) where \(f\)
is non-zero. The existence of the cover \((U_{f_i})\) will be a
consequence of: \oldpage{2-06}

\protect\phantomsection\label{lemma-2}
\begin{itenv}{Lemma 2}
If \(x\) is an arbitrary point of \(V\), and \(U\) is an affine open
subset of \(V\) containing \(x\), then there exists a regular function
\(f\) on \(V\) such that \(V_f=U_f\), and such that \(V_f\) contains
\(x\).

\end{itenv}

\begin{proof}
Let \(W\) be the complement of \(U\) in \(V\), and let
\({\mathscr{I}}(W)\) (resp. \({\mathscr{I}}(W,x)\)) be the sheaf of
germs of functions that are zero on \(W\) (resp. on \(W\) and \(x\)). We
then have an exact sequence: \[
  0 \to {\mathscr{I}}(W,x) \to {\mathscr{I}}(W) \to {\mathscr{I}}(W)|{\mathscr{I}}(W,x) \to 0
\] and the sheaf \({\mathscr{I}}(W)|{\mathscr{I}}(W,x)\) is clearly
non-zero. So there exists a section \(f\) of \({\mathscr{I}}(W)\) that
does not belong to \(\Gamma({\mathscr{I}}(W,x))\); this is the desired
function.
\end{proof}

The quasi-compactness of \(V\), along with \hyperref[lemma-2]{Lemma 2},
imply the existence of an affine cover of \(V\) by a finite number of
\(V_{f_i}\). It remains only to show that the ideal generated by the
\(f_i\) is \(A\). But if \(U\) is an affine open subset, then the
\(U_{f_i}\) cover \(U\), and it follows that restrictions of the \(f_i\)
to \({\mathscr{A}}(U)\) generate the unit ideal. In other words, if
there are \(p\) sections \(f_i\), and if \({\mathscr{A}}^p\) denotes the
direct sum of \(p\) sheaves, each isomorphic to \({\mathscr{A}}\), then
the morphism from \({\mathscr{A}}^p\) to \({\mathscr{A}}\) defined by
the \(f_i\) is surjective. The same is true for the morphism
\(A^p\to A\) defined by the \(f_i\).

\bigskip

For (d)\(\implies\)(a): we will show that the morphism \[
  \varphi\colon (S(V),S({\mathscr{A}})) \to (V(A),{\mathscr{U}})
\] is an isomorphism, and that \(A\) is an algebra of finite type with
no nilpotent elements.

Note first of all that, if \(f\in A\) is a regular function on \(V\),
then the coordinate ring of \(V_f\) is exactly \(A_f\). Indeed, if
\(A_i\) denotes the coordinate ring of \(V_{f_i}\), and
\(A_{ij}=A_{if_j}=A_{jf_i}\) the coordinate ring of
\(V_{f_i}\cap V_{f_j}=V_{f_i\cdot f_j}\), then we have the famous exact
sequence (see Talk no. 1): \[
  0 \to A \to \prod_i A_i \to \prod_{i\neq j}A_{ij}
\] whence (by the exactness of the functor \(M\mapsto M_f\)) we have the
exact sequence: \[
  0 \to A_f \to \prod_i(A_i)_f \to \prod(A_{ij})_f.
\]

But the \((A_i)_f\) and \((A_{ij})_f\) are clearly the coordinate rings
of \(V_{f_i}\cap V_f\) and \(V_{f_i}\cap V_{f_j}\cap V_f\); thus \(A_f\)
is coordinate ring of \(V_f\). \oldpage{2-07}

From this we see that the map \(S(V)\to V(A)\) is injective, and that it
induces an isomorphism \[
  (S(V_{f_i}), {\mathscr{A}}|S(V_{f_i})) \to (V(A)_{f_i}, {\mathscr{U}}|V(A)_{f_i}).
\]

Also, the \(V(A)_{f_i}\) cover \(V(A)\), since the \(f_i\) generate the
unite ideal. By gluing the pieces together, we thus obtain an
isomorphism \[
  (S(V),S({\mathscr{A}})) \to (V(A),{\mathscr{U}})
\] and it remains only to show that \(A\) is of finite type and has no
nilpotent elements.

But \(A\) is an algebra of functions, and has no nilpotent elements.
Now, for all \(i\), let \(a_{i_k}\) be a finite number of elements of
\(A\) such that the images of \(f_i\) and \(a_{i_k}\) in \(A_{f_i}\)
generate the ring \(A_{f_i}\). Finally, let \(a_i\in A\) be such that
\(\sum a_i f_i=1\) in \(A\). We will show that the \(f_i\), \(a_{i_k}\),
and \(a_i\) all together generate the ring \(A\):

Indeed, by taking a suitable power of the formula \(\sum a_i f_i=1\), we
obtain formulas \(\sum a_i(m)f_i^m=1\), where the \(a_i(m)\) can be
expressed in terms of the \(a_i\) and \(f_i\). If now \(x\) is an
element of \(A\), then its image in \(A_{f_i}\) can be written as
\(x_i|f_i^p\) for some \(p\) large enough. It then follows (see Talk no.
1) that, for \(n\) large enough, we have \[
  x = \sum_i a_i(n+p) x_i f_i^n
\] which finishes the proof.
\end{proof}

\protect\phantomsection\label{corollary}
\begin{itenv}{Corollary}
If \(V\) is an affine algebraic set, \(U\) an open subset of \(V\),
\(A\) the coordinate ring of \(V\), and \({\mathscr{U}}(U)\) the
coordinate ring of \(U\), then the following are equivalent:

\begin{enumerate}
\def\labelenumi{\alph{enumi}.}
\tightlist
\item
  \(U\) is an affine open subset.
\item
  If \(M\) is an \(A\)-module, then the canonical homomorphism \[
     M\otimes_A{\mathscr{U}}(U) \to {\mathscr{M}}(U)
   \] is bijective.
\end{enumerate}

\oldpage{2-08}

Furthermore, if one of these equivalent properties is satisfied, then
\({\mathscr{U}}(U)\) is \(A\)-flat.

\end{itenv}

\begin{proof}
(a)\(\implies\)(b): Indeed, \(M\mapsto{\mathscr{M}}(U)\) is then a
right-exact functor in \(M\), where \(M\) runs over all \(A\)-modules.
In particular, if \[
  L_1 \to L_0 \to M \to 0
\] is a resolution of \(M\) by free \(A\)-modules, then we have the
diagram \[
  \begin{CD}
    L_1\otimes_A{\mathscr{U}}(U) @>>> L_0\otimes_A{\mathscr{U}}(U) @>>> M\otimes_A{\mathscr{U}}(U) @>>> 0
  \\@VV{w}V @VV{v}V @VV{u}V @.
  \\{\mathscr{L}}_1(U) @>>> {\mathscr{L}}_0(U) @>>> {\mathscr{M}}(U) @>>> 0
  \end{CD}
\]

Since \(v\) and \(w\) are isomorphisms, so too is \(u\).

\bigskip

(b)\(\implies\)(a): We will show that, if \[
  0 \to {\mathscr{F}}' \to{\mathscr{G}}' \to {\mathscr{H}}' \to 0
\] is an exact sequence of coherent sheaves on \(U\), then the sections
on \(U\) give an exact sequence.

But there exists an exact sequence \[
  0 \to {\mathscr{F}} \to {\mathscr{G}} \to {\mathscr{H}} \to 0
\] of coherent sheaves on \(V\) such that
\({\mathscr{F}}|U={\mathscr{F}}'\), \({\mathscr{G}}|U={\mathscr{G}}'\),
and \({\mathscr{H}}|U={\mathscr{H}}'\) (see \autocite{1}, with the proof
in the appendix of \autocite{2}). So, if \(M\), \(N\), and \(P\) are the
modules associated to \({\mathscr{F}}\), \({\mathscr{G}}\), and
\({\mathscr{H}}\), then the sequence \[
  M\otimes_A{\mathscr{U}}(U) \to N\otimes_A{\mathscr{U}}(U) \to P\otimes_A{\mathscr{U}}(U) \to 0
\] is exact, which proves that the homomorphism \[
  {\mathscr{G}}'(U) = N\otimes_A{\mathscr{U}}(U) \to {\mathscr{H}}'(U) = P\otimes_A{\mathscr{U}}(U)
\] is surjective.
\end{proof}


% Bibliography
\nocite{*}
\printbibliography[heading=bibintoc,title=Bibliography]

\end{document}
